\section{Квантовая механика}
\label{sec:qm}

Классическая траектория не выдерживает атом:
устойчивые орбиты и линейчатый спектр ей противоречат.
Кулоновская задача для электрона,
\begin{equation}
  \Bigl(-\frac{\planckHbar^2}{2\massElectron}\opNabla^2 - \frac{\chargeE^2}{4\symPi\constEpsZero r}\Bigr)\wavePsi = \energyE\wavePsi,
  \label{eq:hydrogen}
\end{equation}
даёт дискретные уровни
\begin{equation}
  E_n = -\frac{\massElectron \speedC^2\alphaFs^2}{2n^2},
  \qquad n=1,2,\ldots
  \label{eq:rydberg}
\end{equation}
Опыт интерференции показывает ещё и то, что складываются не частоты исходов, а амплитуды:
\begin{equation}
  w \propto |\wavePsi_1+\wavePsi_2|^2,
  \label{eq:interference}
\end{equation}
иначе тёмных полос не было бы.
Состояние, из которого такие амплитуды читаются, есть вектор;
частота исхода на макроэтаже есть квадрат его модуля,
\begin{equation}
  w(a) = |\langle a|\wavePsi\rangle|^2.
  \label{eq:born}
\end{equation}
Обратимость, уже видная в \eqref{eq:newton}, здесь означает,
что эволюция не меняет полную норму.
Уравнение Шрёдингера
\begin{equation}
  i\planckHbar\,\opPartialT\wavePsi = \hamiltonianH\wavePsi,
  \qquad \hamiltonianH^\dagger = \hamiltonianH,
  \label{eq:schrodinger}
\end{equation}
сохраняет
\begin{equation}
  \langle\wavePsi|\wavePsi\rangle = \mathrm{const}
  \label{eq:norm}
\end{equation}
и по состоянию в один момент восстанавливает состояние в предыдущий.

Опыт со спином электрона даёт два исхода,
\begin{equation}
  s_z = \pm\frac{\planckHbar}{2},
  \label{eq:spin-z}
\end{equation}
и ещё одно свойство поворота.
Поворот на угол~$\symTheta$ вокруг оси~$\mathbf{n}$ действует как
\begin{equation}
  U(\mathbf{n},\symTheta)
  = \exp\bigl(-i\symTheta\,\mathbf{n}\cdot\boldsymbol{\sigma}/2\bigr).
  \label{eq:spin-rotation}
\end{equation}
Оборот на $2\symPi$ меняет знак, оборот на $4\symPi$ возвращает состояние:
\begin{equation}
  U(2\symPi)=-\mathbb{1},
  \qquad
  U(4\symPi)=+\mathbb{1}.
  \label{eq:spin-4pi}
\end{equation}
Запись состояния --- столбец $\wavePsi\in\spaceC^2$.
Общий множитель-фаза $e^{i\symAlpha}$ не меняет частот \eqref{eq:born}
и соответствует сохранению заряда, уже записанному \eqref{eq:continuity}.
Локальный поворот двух компонент лежит в $SU(2)$.

Физическая величина, построенная из такого состояния,
не имеет права зависеть от того, где проведён разрез аргумента на $2\symPi$.
Комбинация, которая разреза не видит, собирается из амплитуды и сопряжённой амплитуды:
\begin{equation}
  \langle\wavePsi|O|\wavePsi\rangle.
  \label{eq:observable}
\end{equation}

Квантовая механика закрывает эти факты для одной системы с заданным числом степеней свободы.
Одиночный отсчёт прибора по-прежнему стоит рядом с уравнением.
Когда степеней свободы много и частицы рождаются и исчезают,
та же норма и тот же световой конус записываются как теория поля.

